\section*{Введение}
\addcontentsline{toc}{section}{Введение}

\textbf{Актуальность темы.} Современное сельское хозяйство переживает этап
цифровой трансформации, в рамках которой всё большее значение приобретают
системы мониторинга и автоматизированного управления ресурсами.
Выращивание тропических культур, в частности личи (\textit{Litchi chinensis}),
в условиях Гватемалы характеризуется высокой зависимостью урожайности от
режима орошения. Традиционные подходы к управлению ирригацией основаны на
ручном контроле, что приводит к нерациональному расходу водных ресурсов
и снижению качества продукции. При этом внедрение реальных IoT-датчиков
требует значительных капиталовложений, недоступных для малых фермерских
хозяйств. Разработка веб-приложения, сочетающего имитационную модель
агрономических процессов с механиками интерактивного взаимодействия
(по аналогии с виртуальными питомцами) и средствами мониторинга
в реальном времени, является перспективным направлением повышения
вовлечённости операторов и эффективности агротехнических процессов.

\textbf{Степень разработанности проблемы.} Задачи моделирования ирригации
активно исследуются в рамках концепции точного земледелия (Precision
Agriculture). Значительный вклад в область внесли профессиональные
агрономические симуляторы --- AquaCrop (ФАО ООН), DSSAT, APSIM, ---
позволяющие прогнозировать урожайность при различных сценариях полива.
Однако существующие модели ориентированы на учёных-аналитиков, имеют
высокий порог вхождения и полностью лишены элементов геймификации,
способствующих повышению мотивации персонала малых ферм. Развлекательные
симуляторы (Farming Simulator) и потребительские приложения с механикой
Тамагочи (Plantagotchi, Senso) демонстрируют эффективность игровых
механик, но не содержат адекватных математических моделей и не
масштабируются на уровень плантаций.

\textbf{Объект исследования} — процессы управления ирригацией
на плантациях тропических культур.

\textbf{Предмет исследования} — программные средства имитационного
моделирования и интерактивные механики для автоматизации мониторинга
и управления поливом в рамках веб-ориентированной платформы.

\textbf{Целью данной работы} является проектирование и реализация
клиент-серверного веб-приложения, обеспечивающего мониторинг
состояния растений и интерактивное управление ирригацией
на основе серверной имитационной модели с использованием
элементов геймификации.